\begin{abstract}
Implicit Chain-of-Thought (CoT) methods offer a token-efficient alternative to explicit CoT reasoning in Large Language Models (LLMs), but a persistent performance gap has limited their adoption.
We identify a core \textbf{latent instability issue} when scaling the computational budget of implicit CoT: as the number of reasoning tokens increases, training often becomes unstable and collapses.
Our analysis shows that this instability arises from latent representations becoming homogeneous and losing semantic diversity, caused by insufficient step-level supervision in current implicit CoT methods.
To address this, we propose \textbf{\methodname}, a plug-and-play training module that introduces step-level supervision to stabilize and enrich the latent reasoning space.
\methodname employs an auxiliary decoder during training to align each implicit token with its corresponding explicit reasoning step, ensuring latent states capture distinct and meaningful information.
The auxiliary decoder is removed at inference, preserving the efficiency of implicit CoT with no added overhead.
It also provides interpretability by projecting each latent token onto an explicit reasoning vocabulary, enabling per-step visualization and diagnosis.
\methodname significantly improves both in-domain accuracy and out-of-domain stability of implicit CoT methods, boosting Coconut by +8.2\% on GPT-2 and CODI by +3.0\% on LLaMA-3.1 8B.
It further surpasses the explicit CoT baseline on GPT-2 by 2.1\% with 2.3$\times$ greater token efficiency, while closing the performance gap on larger models like LLaMA-3.1 8B.
% Code: \href{https://github.com/InternLM/SIM-CoT}{https://github.com/InternLM/SIM-CoT}.
\end{abstract}

\section{Introduction}
``\textit{Measure what is measurable, and make measurable what is not so.}''~~~~~~~~~~~~~~~~---~~~Galileo Galilei

The strong reasoning capabilities of Large Language Models (LLMs) \citep{gpt4o, gemini, anthropic} are often unlocked through explicit Chain-of-Thought (CoT) prompting \citep{wei2023chainofthoughtpromptingelicitsreasoning}.
The explicit CoT approach enables LLMs to solve complex problems in a step-by-step manner, yielding high performance in domains like mathematics and programming \citep{guo2025deepseek, muennighoff2025s1}.
Despite its advantages, explicit CoT also faces several limitations.
For instance, explicit CoT approaches must verbalize intermediate thoughts from a fixed vocabulary, thereby precluding the exploration of alternative solution paths \citep{li2025implicit,zhang2025soft}.
Additionally, the generation of extensive intermediate sequences significantly increases inference cost and can result in redundant over-thinking steps or unnecessary verbosity \citep{chen2024not}.

To address the flexibility and efficiency issues of explicit CoT methods, recent \textbf{implicit CoT} approaches \citep{hao2024coconut, zhang2025soft, li2025implicit} have been proposed by representing reasoning in a continuous latent space rather than as a sequence of discrete text tokens.
The implicit CoT methods allow each latent representation to encode richer information than a single explicit token, often with a significantly smaller number of latents than the length of an explicit reasoning chain.
Early representative implicit work like Coconut \citep{hao2024coconut} improves efficiency while still capturing useful intermediate structure.
More recent approaches, such as CODI \citep{shen2025codi}, further apply trajectory-level distillation from explicit reasoning paths to enhance performance.
Despite these advancements, a \textbf{performance gap} still exists between existing implicit CoT methods and their explicit counterparts. The implicit CoT approaches are \textit{fast}, \textit{token-efficient} but \textit{less accurate}, which currently limits their broader application.

To narrow the performance gap, inspired by the success of explicit CoT that scales computational budget for better performance, we explore a similar strategy for implicit CoT methods by increasing the number of implicit tokens.
However, in \cref{fig:analysis_all} \textbf{(a)}, we reveal one underlying \textbf{latent instability issue} in current implicit CoT approaches.
As we extend the number of implicit tokens from the default three \citep{hao2024coconut} to five, the training process initially improves accuracy but becomes unstable and sometimes collapses entirely.
To interpret the \textbf{latent instability issue}, we analyze implicit tokens from models trained on math reasoning data GSM8K-Aug \citep{icot}.
We follow previous works \citep{hao2024coconut,icot} to project the implicit tokens through the LM head and examine their top decoded tokens for analysis.
As shown in \cref{fig:analysis_all} \textbf{(b)}, failed models tend to collapse into homogeneous latent states.
While successful reasoning requires capturing both numerical and operator information, the implicit tokens of failed models primarily represent numbers, almost completely losing the critical operator information.
\cref{fig:analysis_all} \textbf{(c)} further demonstrates that a model's collapse is accompanied by two changes: a reduction in the inter-latent distance and a drift of the latent states away from the central vocabulary embedding space.
The latent representations of failed models become too similar and lose their semantic connection to the tokens they are meant to represent.
\cref{fig:analysis_all} \textbf{(d)} provides an example of the semantic homogenization.
A normal model (top) maintains a large distance between its two latent tokens, allowing them to capture distinct information for numbers and operators.
In contrast, a failed model (bottom)'s latent tokens become homogeneous, with both states decoding to similar information, primarily numbers.

\input{figures/analysis_all}

Our observation (\cref{fig:analysis_all}) reveals the reasons for the latent instability issue: a lack of sufficient step-level supervision for existing implicit methods to maintain the rich and varied internal representations.
Without stronger guidance, the latent space collapses, losing its diversity and making it impossible to reliably encode the distinct, step-level reasoning needed for complex reasoning tasks.
Motivated by our findings, we propose \textbf{S}upervised \textbf{IM}plicit-CoT (\textbf{\methodname}), a plug-and-play module that introduces step-level supervision for implicit CoT approaches to alleviate the latent instability issue.
Instead of supervising only the final answer \citep{hao2024coconut} or the trajectory \citep{shen2025codi}, \methodname uses an auxiliary decoder to align each implicit token with its corresponding explicit reasoning step during training.
The step-level supervision for implicit tokens stabilizes optimization, prevents collapse, and ensures that latent tokens capture meaningful reasoning content.
Crucially, because the auxiliary decoder is removed during inference, our approach incurs virtually no extra computational cost, making it as efficient as standard implicit CoT approaches.
Beyond \textit{accuracy}, \textit{stability}, and \textit{efficiency}, the auxiliary decoder also affords \textit{interpretability} of implicit reasoning.
During training, it defines a projection from latent tokens to the explicit reasoning vocabulary, enabling us to decode each latent step into a human-interpretable summary for verification or error diagnosis.

Experiments show that \methodname acts as a plug-and-play module that boosts both accuracy and stability.
We show that \methodname can be effortlessly combined with various implicit CoT approaches such as Coconut \citep{hao2024coconut}, CODI \citep{shen2025codi}, and training-free approaches \citep{zhang2025soft} to further enhance reasoning performance.
On GPT-2, \methodname surpasses both the strong explicit baseline (supervised fine-tuning on explicit CoT data) by 2.1\%, and outperforms existing implicit methods Coconut and CODI by 8.2\% and 4.3\%, respectively.
The performance trend holds as the method scales to larger models such as the LLaMA series. $\text{\methodname}$ achieves improvements over CODI of 3.4\% (LLaMA-3.2 1B), 1.5\% (LLaMA-3.2 3B), and 3.0\% (LLaMA-3.1 8B), in addition to a 9.0\% gain over Coconut on the LLaMA-3.2 1B model.
Furthermore, while previous implicit CoT approaches (e.g., Coconut) collapse when scaled to 8 or 16 implicit tokens, \methodname remains stable and continues to boost performance.

In summary, our contributions are as follows:  
\textbf{1)} We provide a systematic analysis of the latent instability issue of implicit CoT approaches, showing that instability and collapse arise from insufficient supervision.
\textbf{2)} We introduce \methodname, which applies step-level supervision to the model's implicit tokens. \methodname not only integrates seamlessly with existing implicit CoT approaches and boosts performance with minimal inference overhead, but also affords interpretability of implicit reasoning by projecting each latent token onto an explicit reasoning vocabulary, enabling per-step visualization of semantic roles and diagnosis.
\textbf{3)} Through extensive experiments, we demonstrate that \methodname not only improves accuracy in the in-domain dataset, but also generalizes effectively to out-of-domain datasets. The performance gains are consistent across a range of LLMs, including GPT-2 and recent LLaMA 3 models (1B, 3B, and 8B).

\appendix
\newpage
\section*{\centering APPENDIX}

\section*{Usage of Large Language Models}
In this paper, we used LLMs only for minor language polishing and formatting, without generating ideas, analyses, or experimental results.

\section*{Outline}
In this appendix, we provide additional analyses and supporting materials to complement the main text.  
First, in Sec.~\ref{appendix:scaling_backbones}, we investigate the scaling behavior of \methodname on larger backbones and show its robustness across different model sizes.  
In Sec.~\ref{appendix:decoder_size}, we analyze the impact of decoder size on performance, highlighting the trade-offs between moderate and excessively large decoders.  
In Sec.~\ref{appendix:softthinking}, we present experiments on combining soft thinking with \methodname, including setup, results, and a detailed formulation with pseudocode.  
In Sec.~\ref{appendix:related}, we provide an overview of related work in explicit and implicit chain-of-thought reasoning.  
In Secs.~\ref{appendix:implementation} and \ref{app:train-infer-details}, we describe our implementation, hyperparameter configurations, and training/inference procedures, including benchmark and dataset details.  
In Sec.~\ref{appendix:simcot}, we introduce the SIM-CoT training procedure and provide pseudocode for step-level supervision.  
In Sec.~\ref{appendix:geometry}, we offer geometric diagnostics of the latent space, analyzing inter-latent distances and distance to the vocabulary center.  
In Sec.~\ref{appendix:interpretability}, we discuss interpretability analysis, including latent visualization, distance analysis, and summary findings.  
Finally, we provide future directions, declarations on LLM usage and additional case studies on GSM8k to further illustrate the reasoning process and visualization choices.


\section{Additional Analysis on Scaling to Larger Backbones}
\label{appendix:scaling_backbones}

As shown in Table~\ref{tab:llama3-3b_and_8b}, \textbf{On LLaMA 3.2 3B}, \methodname improves over CODI by $+1.5$ points on GSM8k-Aug and $+1.6$ points on SVAMP, while maintaining comparable performance on GSM-Hard and MultiArith.  
This indicates that step-level implicit supervision can strengthen strong implicit reasoning baselines.  

\textbf{On LLaMA 3.1 8B}, \methodname yields $+3.0$ points on GSM8k-Aug, $+1.3$ points on SVAMP, and $+0.8$ points on MultiArith relative to CODI, with stable accuracy on GSM-Hard.  
Compared with SFT-CoT, it achieves higher accuracy on MultiArith ($100.0$ vs.\ $98.3$) and SVAMP ($79.4$ vs.\ $73.1$), while remaining similar on GSM-Hard.  

These results confirm that \methodname scales effectively to larger backbones, providing consistent gains while reducing reliance on trajectory-level supervision.

\section{Additional Analysis on Decoder Sizes}
\label{appendix:decoder_size}

Compared to the baseline, as shown in Table~\ref{tab:llama1b-vs-distance}(a), integrating a 1B-scale decoder leads to consistent improvements across all benchmarks.  
However, simply replacing the decoder with larger versions (3B or 8B) does not bring further benefits and slightly degrades performance.


These results suggest that moderate decoder scaling can enhance reasoning ability, but excessively large decoders may introduce optimization difficulties or misalignment with the backbone, ultimately limiting generalization.  
Another plausible explanation is that the 1B encoder and 1B decoder originate from the same model, thereby sharing a more compatible representation space that facilitates learning.  
In contrast, larger decoders (3B or 8B) may require additional projection to align with the 1B backbone, which could introduce mismatches and hinder training stability.


\section{Additional Analysis on Soft Thinking}
\label{appendix:softthinking}

Soft thinking \citep{zhang2025soft, wu2025llms} is a training-free method for implicit reasoning in which the latent space is represented as a weighted average over the vocabulary embedding space.  
In contrast, \methodname learns latent representations directly from data during training.  
To our knowledge, no prior work has evaluated a combination of these two approaches; our experiments provide the first such evaluation.

\subsection{Setup}
We apply the proposed soft thinking mechanism on top of both Coconut and SIM-CoT, while adopting GPT-2 as the backbone model.
To assess the effectiveness of this approach, we perform evaluations on a diverse set of mathematical reasoning benchmarks.
The in-domain evaluation is carried out on GSM8k-Aug, which provides augmented training and testing samples closely aligned with the original GSM8k distribution.
To further examine generalization beyond the training domain, we include three out-of-domain benchmarks: GSM-Hard, which contains more challenging arithmetic problems with subtle variations in reasoning steps; MultiArith, which evaluates performance on multi-step arithmetic operations requiring careful sequencing of addition, subtraction, multiplication, and division; and SVAMP, which focuses on variations of elementary word problems designed to test robustness to superficial changes in problem statements.


\subsection{Results}
Table~\ref{tab:soft-thinking} (b) reports the results.  
Adding soft thinking improves accuracy in most cases.  
For Coconut, improvements are observed on GSM-Hard (+0.2) and MultiArith (+1.7), with a slight decrease on SVAMP (-0.2).  
For SIM-CoT, soft thinking consistently enhances performance: GSM8k-Aug (+0.2), GSM-Hard (+0.1), MultiArith (+0.7), and SVAMP (+0.1).

\subsection{Formulation}
Let $z \in \mathbb{R}^d$ denote a continuous latent token, and $E \in \mathbb{R}^{|\mathcal{V}| \times d}$ be the embedding matrix of the vocabulary $\mathcal{V}$.  
Our goal is to enrich the representational capacity of $z$ by incorporating soft thinking, which allows the latent space to draw information not only from its continuous representation but also from the semantic structure of the vocabulary. The process can be described in three steps.

\paragraph{Step 1. Vocabulary distribution.}  
The continuous latent token $z$ is first mapped into a probability distribution over the vocabulary space:
\[
p = \text{softmax}(W z),
\]
where $W \in \mathbb{R}^{|\mathcal{V}| \times d}$ is the output projection matrix and $p \in \mathbb{R}^{|\mathcal{V}|}$ is the resulting distribution.  
This step can be viewed as interpreting the latent token in terms of vocabulary-level semantics, where each token in $\mathcal{V}$ is assigned a likelihood according to its relevance to $z$.

\paragraph{Step 2. Soft-thinking embedding.}  
Using the distribution $p$, we compute a weighted mixture of vocabulary embeddings:
\[
z_{\text{soft}} = E^\top p = \sum_{v \in \mathcal{V}} p_v \, E_v ,
\]
where $E_v$ is the embedding vector corresponding to token $v$.  
This operation can be seen as constructing a "soft token" that captures multiple semantic hypotheses simultaneously, instead of committing to a single discrete vocabulary token. As a result, $z_{\text{soft}}$ provides richer and smoother information than a hard token lookup.

\paragraph{Step 3. Combination.}  
Finally, we combine the original continuous latent $z$ with the soft-thinking embedding $z_{\text{soft}}$:
\[
z' = \alpha z + \beta z_{\text{soft}},
\]
where $\alpha = \texttt{continuous\_weight}$ and $\beta = \texttt{soft\_weight}$ are hyperparameters that balance the contribution of the continuous and soft-thinking components.  
This formulation allows $z'$ to retain the model’s learned continuous representations while also grounding them in the vocabulary space. Intuitively, the continuous part encourages compact reasoning within the latent space, whereas the soft-thinking component brings in semantic priors from the vocabulary, leading to more stable and interpretable reasoning.  

The pseudocode implementation of the above process is presented as follows.

\begin{algorithm}[h]
\caption{Soft Thinking with Continuous Tokens}
\label{alg:soft-thinking}
\begin{algorithmic}[1]
\Require Continuous latent $z$, embedding matrix $E$, weights $\alpha, \beta$
\If{$\beta > 0$}
    \State Compute logits: $l \gets Wz$
    \State Convert to probabilities: $p \gets \text{softmax}(l)$
    \State Form soft embedding: $z_{\text{soft}} \gets E^\top p$
    \State Update latent: $z' \gets \alpha z + \beta z_{\text{soft}}$
\Else
    \State $z' \gets z$
\EndIf
\State \Return $z'$
\end{algorithmic}
\end{algorithm}


\subsection{Analysis}
The results demonstrate that soft thinking complements training-based implicit reasoning.
The hybrid latent $z'$ integrates semantics learned through training and distributional information from vocabulary mixing, which enables the model to explore diverse intermediate states rather than committing to a single deterministic path.
This leads to improvements in both in-domain and out-of-domain benchmarks.  
Our findings suggest that combining training-free construction with training-based supervision provides gains beyond either approach in isolation.

\input{tables/scale_decoder_and_distance_analysis}
\input{tables/ablation_soft_thinking}

\section{Related Work}
\label{appendix:related}

\noindent \textbf{Explicit chain-of-thought reasoning.}  
Chain-of-thought (CoT) prompting enables large language models (LLMs) to generate intermediate reasoning steps before producing the final answer \citep{wei2023chainofthoughtpromptingelicitsreasoning}.  
This approach has been widely studied and extended in many directions. Self-consistency samples multiple reasoning paths and selects the majority answer to improve reliability \citep{wang2023selfconsistency}.  
Least-to-most prompting decomposes a complex question into simpler sub-problems and solves them in order \citep{zhou2023least}.  
Reflection-based reasoning allows the model to revise or verify its own intermediate steps, leading to better correctness \citep{shinn2023reflexion, madaan2023self}.  
Other works focus on using external tools or symbolic solvers together with explicit reasoning, which further improves accuracy in mathematics and program synthesis \citep{yao2023react}.  
Methods such as progressive-hint prompting \citep{zheng2023progressive} and step-level feedback \citep{wei2025survey} study how supervision can be incorporated into explicit reasoning to make reasoning more structured.  
Despite these advances, explicit CoT has clear drawbacks. Because it generates long token sequences, inference cost grows rapidly with reasoning length, and many intermediate steps are redundant or irrelevant to the final answer.  
Moreover, since explicit reasoning is restricted to tokens from a fixed vocabulary, it often commits to a single trajectory and shows limited reasoning diversity \citep{li2025implicit, zhang2025soft, xu2025softcot}. 

\noindent \textbf{Implicit chain-of-thought reasoning.}
Implicit CoT performs multi-step computation in a continuous latent space instead of emitting long textual traces, reducing decoded length while keeping internal structure. Prior work follows four practical routes. First, \textbf{knowledge internalization} trains models to carry out reasoning internally by progressively removing explicit traces or by using dedicated control embeddings; examples include iCoT-SI \citep{icot}, which removes steps during training to internalize reasoning. Second, \textbf{architectural modification} controls compute by reusing or skipping layers, or by adding light recurrence, so models can refine hidden states without lengthening outputs \citep{loop,innerthinking,ccot,tokenassorted,cotformer,geiping2025scaling}. Third, \textbf{training-free} methods construct continuous latents directly from the model’s probability distribution over the vocabulary; Soft Thinking mixes embeddings by probability to form “concept” tokens that explore alternative paths without updating weights, which improves efficiency and diversity but does not bind each latent to step-level semantics \citep{zhang2025soft,wu2025llms}. 

The fourth route, \textbf{auto-regressive latent reasoning}, updates and concatenates latent states in place of some token-level decoding and is the most relevant to our work \citep{xu2025softcot,tan2025colar}. Coconut applies \textbf{answer-level} supervision—training on the final answer while leaving intermediate latents weakly constrained \citep{hao2024coconut}. CODI adds \textbf{trajectory-level} distillation by aligning an implicit trajectory with an explicit CoT trace, narrowing the gap to explicit CoT but giving only coarse guidance to intermediate steps \citep{shen2025codi}. However, the implicit token length in CODI is fixed during training, which limits its flexibility and makes it less suitable for scaling to variable or longer reasoning chains.
Our framework remains in the auto-regressive setting but changes the supervision: during training, each latent is aligned with its corresponding textual step (\textbf{step-level} supervision), distributing learning signals across the full latent chain to improve stability and semantic fidelity of intermediate states; at inference, the decoder is discarded, ensuring that the decoding cost remains identical to that of standard implicit CoT methods (e.g., Coconut).

\section{Implementation and Training Details}
\label{appendix:implementation}

We provide the full hyperparameter settings, training procedures, and additional analysis used in our experiments. Unless otherwise specified, we use the Adam optimizer with $\beta_1=0.9$, $\beta_2=0.999$, and weight decay of $0.1$. Batch size is set to 128 for GPT-2 and LLaMA 1B, and 64 for LLaMA 3B and 8B. Early stopping is applied with a patience of 3 epochs. We now describe the training setups for Coconut, CODI, and our \methodname.

\subsection{Coconut Training Setup}
Following \citet{hao2024coconut}, GPT-2 and LLaMA 1B \citep{Radford2019LanguageMA,llama3.2} are trained with a fixed learning rate of $1 \times 10^{-4}$. One implicit latent corresponds to two implicit tokens. A curriculum is applied: every three epochs, one explicit reasoning step is replaced by an implicit latent until the maximum number of latent steps is reached. After this expansion, training continues for 15 additional epochs.

\subsection{CODI Training Setup}
For larger backbones such as LLaMA 3B and LLaMA 8B, we adopt task-specific hyperparameter settings to ensure stable training. In particular, we use a learning rate of $3 \times 10^{-4}$ for LLaMA 3B and train for 8 epochs, while for LLaMA 8B the learning rate is reduced to $1 \times 10^{-4}$ with 6 training epochs. These choices are motivated by the increased sensitivity of larger models to optimization dynamics, where smaller learning rates and fewer epochs help to prevent overfitting and instability.  

When reproducing CODI on GPT-2 and LLaMA 1B, we strictly follow the configurations reported by \citet{shen2025codi}. Specifically, we use a learning rate of $3 \times 10^{-3}$ with 40 epochs for GPT-2, and a learning rate of $8 \times 10^{-4}$ with 10 epochs for LLaMA 1B. Adopting these settings ensures that our results are directly comparable to prior work and isolates the effect of our proposed method, rather than confounding it with differences in optimization schedules.  


\subsection{Summary of Hyperparameters}
\begin{table}[h!]
\centering
\caption{Training hyperparameters across different models.}
\label{tab:hyperparams}
% \large
\begin{tabular}{lccc}
\toprule
\textbf{Model} & \textbf{Method} & \textbf{LR} & \textbf{Epochs} \\
\midrule
GPT-2          & Coconut & $1 \times 10^{-4}$ & 15 + curriculum  \\
LLaMA 1B       & Coconut & $1 \times 10^{-4}$ & 15 + curriculum  \\
GPT-2          & CODI    & $3 \times 10^{-3}$ & 40  \\
LLaMA 1B       & CODI    & $8 \times 10^{-4}$ & 10  \\
LLaMA 3B       & CODI    & $3 \times 10^{-4}$ & 8   \\
LLaMA 8B       & CODI    & $1 \times 10^{-4}$ & 6   \\
\bottomrule
\end{tabular}
\end{table}

\section{Training and Inference Details}
\label{app:train-infer-details}

\paragraph{Curriculum for \(K\).}  
We use a curriculum schedule to gradually increase the number of implicit steps. Each latent corresponds to two implicit tokens. Let \(K_{\max}\) denote the maximum number of latents. Starting from \(K^{(0)}=0\), the number of implicit steps after epoch \(e\) is
\[
K^{(e)} \;=\; \min\!\left(K_{\max},\;\Big\lfloor \tfrac{e}{\Delta e} \Big\rfloor \right),
\]
where \(\Delta e\) is the update interval in epochs. Once \(K^{(e)}\) reaches \(K_{\max}\), it remains fixed for the remainder of training.

\paragraph{Inference and Efficiency.}  
At inference time, the auxiliary decoder is removed and only the base model is executed:
\[
U^{(0)} = [\, e(x_1), \ldots, e(x_T) \,], \quad 
\text{for } k=1,\ldots,K:\;
z_k = H_\theta(U^{(k-1)}), \quad
U^{(k)} = U^{(k-1)} \oplus z_k ,
\]
and after $K$ implicit steps the model switches back to explicit decoding to generate the final answer sequence $a$ as in Eq.~\eqref{eq:explicit}.  
The total decoding length is $T + K + L_a$, where $T$ is the input length, $K$ the number of implicit steps, and $L_a$ the answer length. In practice, the cost is comparable to other implicit reasoning methods because $K$ is moderate. In tasks where explicit CoT requires long trajectories ($L_{\text{CoT}} \gg K$), the implicit formulation reduces decoding positions, providing efficiency gains without loss of reasoning accuracy.

\paragraph{Benchmark Detail.}
\label{app:benchmark_detail}
\textbf{SVAMP} \citep{patel-etal-2021-nlp} (Simple Variations on Arithmetic Math Word Problems) is a benchmark dataset designed to test the robustness of math word problem solvers to superficial changes. It contains 1,000 elementary-level arithmetic word problems (grade 4 and below), each involving a single unknown and solvable by an arithmetic expression with no more than two operators. The problems are transformations of existing datasets (such as MAWPS and ASDiv-A) with controlled variations in wording, structure, and number values to reduce artifacts and superficial cues. SVAMP’s average number of reasoning steps required is around 1.2, similar to the base datasets, but model performance drops significantly when tested on SVAMP, showing that many models rely on heuristic patterns rather than deep understanding. 

\textbf{GSM-Hard}~\citep{gao2022pal} is a more challenging variant of the GSM8K dataset, intended to test models’ ability to cope with harder numerical values. It retains the same problem statements as the original GSM8K, but replaces many of the numbers with larger and less common numerical values, making superficial arithmetic computation and reasoning harder. The dataset has about 1,319 examples (matching the GSM8K test set size).

\textbf{MultiArith}~\citep{roy2016solvinggeneralarithmeticword} is a dataset of multi-step arithmetic word problems designed to challenge systems to correctly sequence multiple operations. It contains 600 problems collected from educational sources, each solvable by one equation involving two or more of the four basic operations (addition, subtraction, multiplication, division). The problems require reasoning over multiple sentences to extract and combine numeric quantities, understand the implied operations, and compute the result. MultiArith has been widely used as a benchmark for evaluating arithmetic reasoning generalization, particularly for models that go beyond single-operation problems. Analyses show that many models struggle on these examples compared to simpler datasets, highlighting the importance of handling compositional and sequential numerical reasoning.

\paragraph{Training Data.}
\label{app:training_data}
\input{figures/gsm8k_step_distribution}
\textbf{GSM8K-Aug}~\citep{icot} is the only training corpus we use. It is an augmented dataset derived from GSM8K~\citep{cobbe2021trainingverifierssolvemath}, expanding the original 8.5k training problems to roughly 385k examples through paraphrasing, numerical resampling, and synthetic generation with GPT-4. The distribution of reasoning steps in GSM8K-Aug is illustrated in Fig.~\ref{fig:gsm8k-step-distribution}, where the majority of problems require two to four steps, while a long-tail of six or more steps persists. This balance of common and complex instances makes GSM8K-Aug particularly suitable for training models that need to generalize across reasoning difficulty levels.

\section{SIM-CoT Training Implementation}
\label{appendix:simcot}

We provide pseudocode for the SIM-CoT training process, which illustrates how continuous latent embeddings are aligned with explicit supervision at the step level. In particular, each reasoning step in the explicit chain is mapped to a corresponding latent representation, and the training objective enforces consistency between the predicted latent tokens and the ground-truth step annotations. This design ensures that the model learns to represent intermediate reasoning steps in a compact latent space while still retaining interpretability through explicit alignment. By supervising at the step level rather than only at the final answer or trajectory level, SIM-CoT enables finer control over the reasoning process and reduces instability that often arises when scaling to longer chains or larger numbers of implicit tokens.


\begin{algorithm}[h]
\caption{SIM-CoT Training Procedure}
\label{alg:simcot}
\begin{algorithmic}[1]
\Require Batch size $b$, number of thoughts $C$, continuous embeddings $Z$, tokenized inputs $X$, embedding matrix $E$
\For{each thought $t = 1, \ldots, C$}
    \For{each sample $i = 1, \ldots, b$}
        \State Extract continuous embeddings $z_{i,t}$ from $Z$
        \State Obtain token embeddings $e_{i,t}$ from $E(X_{i,t})$
        \State Concatenate embeddings: $h_{i,t} \gets [z_{i,t}; e_{i,t}]$
        \State Build attention mask $m_{i,t}$ up to EOS
        \State Assign position ids $p_{i,t}$
        \State Prepare labels $y_{i,t}$ with masked tokens set to $-100$
    \EndFor
\EndFor
\State Pad and stack $\{h,m,p,y\}$ to maximum sequence length
\State Prepare 4D attention mask: $\hat{M} \gets \text{PrepareMask}(M)$
\State Forward pass: $\hat{O} \gets \text{ExplainableLLM}(H, \hat{M}, P)$
\State Extract logits: $L \gets \hat{O}.\text{logits}$
\State Shift logits and labels: $L' \gets L[:, :-1], \; Y' \gets Y[:, 1:]$
\State Compute cross-entropy loss: $\ell = \text{CrossEntropy}(L', Y')$
\State Normalize $\ell$ over valid positions
\Ensure Final training loss $\ell$
\end{algorithmic}
\end{algorithm}

\section{Geometric Diagnostics of the Latent Space}
\label{appendix:geometry}

We analyze the geometry of latent representations with two metrics.

\textit{Inter-latent distance.}  
\begin{equation}
\mathrm{Dist}(z_{1:K})
= \frac{2}{K(K-1)} \sum_{1 \le i < j \le K} \big\| z_i - z_j \big\|_2 .
\end{equation}
A larger value indicates better separation, reducing the risk of collapse.

\textit{Distance to vocabulary center.}  
Let $\mu = \tfrac{1}{|\mathcal{V}|}\sum_{v\in\mathcal{V}} E_v$ denote the mean embedding. Then,
\begin{equation}
\mathrm{DistVC}(z_{1:K})
= \frac{1}{K} \sum_{k=1}^{K} \big\| z_k - \mu \big\|_2 .
\end{equation}
Moderate values indicate that latents remain close enough to the lexical manifold for stability, while avoiding collapse toward the center. These diagnostics are not used in training but serve as indicators of diversity and stability in the learned latent space.

\input{figures/distance_visulize}

\section{Additional Details for Interpretability Analysis}
\label{appendix:interpretability}

\subsection{Making Implicit Reasoning Visible}
Continuous thoughts produced by implicit reasoning models are represented as latent embeddings that do not correspond to discrete vocabulary tokens, and therefore cannot be directly decoded by a tokenizer. This makes it difficult to interpret how the model internally organizes multi-step reasoning. To address this, we reuse the decoder that was employed for step-level supervision during training, and apply it at inference time to map each latent embedding into a human-readable token sequence.  

As illustrated in Figure~\ref{fig:visualization}, the process begins with a natural language problem (e.g., a math word problem) that is embedded and passed into the large language model. The model generates a sequence of implicit latent tokens, which capture intermediate reasoning steps in continuous space. These latent tokens are then fed into the optional decoder, which translates them into interpretable expressions. Each latent corresponds to one reasoning step, and the autoregressive generation order encodes the dependency structure across steps.  

For example, in the GSM8k case study shown in the figure, the first latent is decoded as $0.3 \times 120 = 36$, representing the number of watermelons harvested initially. The second latent builds upon this result to compute $120 - 36 = 84$, the remaining melons. The third latent then calculates $\tfrac{3}{4} \times 84 = 63$, and the final latent derives the answer $84 - 63 = 21$. For clarity and to save space, the figure merges the second and third steps into a single box, but the actual implicit reasoning unfolds across four distinct latent steps. This sequence of decoded latents mirrors the logic of explicit chain-of-thought reasoning, while being produced implicitly within the latent space.
  

By projecting implicit tokens into interpretable space, we gain direct visibility into how the model structures multi-step reasoning. This not only enables analysis of the correctness and consistency of intermediate steps, but also highlights the dependencies across steps that underlie the final prediction. The visualization confirms that SIM-CoT can encode semantically meaningful and logically ordered reasoning steps in its latent space, bridging the gap between implicit and explicit reasoning.


\subsection{Distance Analysis}
We analyze the geometry of latent representations by measuring two quantities: the average pairwise distance among latent tokens (Dist.) and the average distance from latent tokens to the vocabulary center (Dist. to VC).  
Table~\ref{tab:llama1b-vs-distance}(b) reports the numerical outcomes under different configurations, while Figure~\ref{fig:distance_visualize} provides a visualization of the same phenomena.

As the number of latent tokens increases from 1 to 5, Dist. gradually grows from 20.30 to 28.34, suggesting that the model distributes the latent representations more sparsely in the embedding space, which improves separability. However, in the failed case with 5 latents, Dist. collapses to 4.21, reflecting a degeneration where all latent tokens converge to nearly identical points, losing their ability to represent distinct reasoning steps. By contrast, SIM-CoT pushes Dist. to 32.81 under the same configuration, showing that step-level supervision effectively enforces stronger separation and prevents collapse.

For Dist. to VC, we observe the highest value with 1 latent (36.20). As the number of latents increases, Dist. to VC decreases toward approximately 28, meaning that the representations become better aligned with the vocabulary manifold rather than drifting away. The failed 5-latent case instead shows an abnormal increase to 39.39, a clear signal of instability where the model drifts away from the vocabulary space. SIM-CoT stabilizes this measure at 29.80, achieving a balance: the latent tokens remain distinct from each other while staying anchored close enough to the vocabulary space to preserve semantic interpretability.

The visualization in Figure~\ref{fig:distance_visualize} corroborates these trends. Normal implicit tokens (first row) are well separated and positioned near the vocabulary manifold. Failed tokens (second row) collapse together and drift outward, losing both separation and grounding. After applying SIM-CoT (third row), the tokens regain a structured configuration: distances between latents expand, while their proximity to the vocabulary center remains stable. This alignment between quantitative and qualitative results highlights the effectiveness of SIM-CoT in maintaining a structured and interpretable latent space.

\subsection{Summary}
Overall, the results demonstrate that SIM-CoT establishes a balance between \textbf{diversity} and \textbf{stability} in the latent space. Larger inter-latent distances mitigate representation collapse, while moderate distances to the vocabulary center prevent excessive drift. This equilibrium supports stable implicit reasoning and provides robustness when scaling to more latent tokens.

\section{Future Directions}

We outline several promising directions for extending this work:  
\textbf{(1) Multimodal extension.} Incorporating intermediate supervision from images \citep{chen2023sharegpt4v, sun2024alpha, liu2024mmdu} or videos \citep{chen2024sharegpt4video, wei2025videorope} could generalize implicit CoT to multimodal reasoning tasks \citep{shen2025efficient, zhang2025booststep}, an especially timely direction given the rapid progress of multimodal learning \citep{liu2023llava, chen2023sharegpt4v, zhang2025sec, liu2025songgen, wei2025videoropepp, zhang2024internlm, dong2024internlm, zhao2025omnialign, ding2025mmifenginemultimodalinstructionfollowing, chen2024open, xing2025scalecap, xing2024pyramiddrop}. 
\textbf{(2) Multi-path implicit reasoning.} Exploring branched latent structures, inspired by Tree-of-Thought \citep{yao2023tree, wang2023selfconsistency, zhang2023evaluating} methods, may enhance the diversity and robustness of implicit reasoning.  
\textbf{(3) Integration with RLHF.} Step-level supervision can complement preference optimization and reinforcement learning \citep{guo2025deepseek, rafailov2023direct, liu2025visual, internlmxcomposer2_5_reward}, leading to a more reliable and adaptive reasoning framework.  
\textbf{(4) Theoretical foundations.} Future work may also study implicit reasoning from the perspectives of information theory and representation learning \citep{zhu2025reasoning, xu2025cot}, providing a principled explanation for why step-level supervision stabilizes the latent space.



\section{Additional Case Studies on GSM8k}
In practice, implicit reasoning continues to produce latent tokens even after the correct answer has been reached. These trailing latents no longer introduce new steps but simply repeat the final prediction. For clarity, we omit such redundant tokens in the visualizations. As a result, only the latents that correspond to meaningful intermediate steps are displayed in Figure~\ref{fig:case_all}, while those mapping directly to the final answer are hidden. This choice improves readability without changing the underlying reasoning process.

Notably, the decoded reasoning steps consistently match the semantic structure of explicit chain-of-thought annotations, while being generated implicitly within the latent space. The final predictions align with the ground-truth answers, demonstrating that \methodname is capable of encoding interpretable and step-ordered reasoning without requiring explicit supervision at inference time.  

\input{figures/case_all}


\section{Analysis of Implicit COT: the Latent Instability Issue}
We first present an analysis (\cref{fig:analysis_all}) of the limitations in implicit latent CoT approaches.
We follow Coconut \citep{hao2024coconut} and analyze implicit latents by projecting them through the LM head and examining the top-8 decoded tokens to understand the semantic and geometric properties.

\noindent \textbf{Latent Instability Issue.}
\cref{fig:analysis_all} \textbf{(a)} shows the training process of Coconut when the number of implicit latent tokens is progressively increased. 
Initially, as the number of latents increases from one to four, the model's accuracy generally improves, suggesting that using more latents can enhance performance.
However, a significant drop in accuracy occurs when the number of latents is scaled to five, with performance collapsing to its worst point of 12.5\%. 
The latent instability issue indicates that the implicit reasoning approach is sensitive to the choice of the number of latent tokens, as shown by the sharp drop and subsequent fluctuations in accuracy after adding the fifth latent.

\noindent \textbf{Information Loss.}
\cref{fig:analysis_all} \textbf{(b)} presents an analysis of how different levels of accuracy are affected by the number of latent tokens, using accuracy metrics at three levels: number, operator, and answer.
The bar chart reveals a clear trend: as the number of latent tokens increases from 1 to 5, there is a general decline in performance across all three metrics, especially for the operator accuracy.
The strong correlation between increased latent tokens and declining performance, particularly the sharp fall during failure, suggests that implicit latents do not consistently capture the necessary compositional reasoning process without more explicit, fine-grained supervision.

\noindent \textbf{Shifted Distance.}
\cref{fig:analysis_all} \textbf{(c)} examines the geometric properties of the latent representations during training.
Two metrics are analyzed: the Latent Distance (red), which measures the average distance between pairs of latent vectors, and the Vocab Distance (blue), which measures the average distance from each latent vector to the center of the vocabulary embedding space.
When the latent CoT model collapses, the latent distance decreases sharply, indicating that the latent vectors are collapsing and becoming nearly identical, losing their distinctiveness.
Simultaneously, the vocab distance increases, showing that these collapsing latents are drifting away from the main lexical embedding space and are no longer grounded in the fundamental token representations used by the model.

\noindent \textbf{Semantic Homogenization.}  
\cref{fig:analysis_all} \textbf{(d)} provides a qualitative analysis of the content of the latent tokens in a normal case versus a failed model.
In the normal implicit model (middle), the decoded tokens from the latents are diverse and meaningful.
In the failed implicit model (bottom), the semantic content of the latents becomes highly homogeneous.
Latent 1 and Latent 2 contain mainly numbers, lacking operators or symbolic information needed for calculation. This shows that successful training produces latents with step-wise reasoning, while without explicit supervision, the latent space collapses into uniform numerical forms.

\noindent \textbf{Summary.}
Our analysis across \cref{fig:analysis_all} (a-d) highlights a crucial trade-off between diversity and stability.
When the model collapses, it loses both its diversity (as the latents become too similar) and its stability (as the latents move away from the token space), leading to catastrophic information loss and a complete failure of the reasoning process, as shown by the sharp drop in overall accuracy.
These combined findings show that without proper guidance, the latent space degenerates, losing its ability to represent distinct reasoning steps.
These challenges motivate our proposed method, which introduces \textbf{step-level implicit supervision} to stabilize the training process and enrich unique semantic content of each latent, all while maintaining efficiency during inference.

\section{Methodology}
\input{figures/teaser}

\noindent \textbf{Overview.}
As shown in \cref{fig:coconut_teaser}, early implicit reasoning studies differ mainly in supervision granularity:
\textbf{Coconut} (top left) uses answer-level supervision, while \textbf{CODI} (top right) introduces trajectory-level signals via distillation.
Both remain coarse and do not tell the model which latent should encode which step.
We propose \textbf{\methodname}, which provides \textbf{step-level implicit supervision}:
During an \textbf{implicit phase}, the LLM runs for a fixed number \(K\) of reasoning steps; at each step \(k\) it takes the \textbf{last hidden state} as the implicit latent \(z_k\) and appends it to the sequence as the next “token” vector.
After \(K\) steps, the model switches back to \textbf{explicit} decoding over the vocabulary to generate the final answer.
A decoder is used only in training to align each \(z_k\) with the textual content of the \(k\)-th reasoning step; at inference, the decoder is removed, so the runtime is essentially that of direct answer generation plus \(K\) forward positions, which is far shorter than explicit CoT token lengths.

\subsection{Notation}
Let \(\mathcal{V}\) be the vocabulary and \(E\in\mathbb{R}^{|\mathcal{V}|\times d}\) the token embedding matrix.
A question is \(x=(x_1,\ldots,x_T)\in\mathcal{V}^T\) with embedded prefix
\[
U^{(0)}=\big(e(x_1),\ldots,e(x_T)\big), \quad e(\cdot)\in\mathbb{R}^{d}.
\]
We run an autoregressive LLM \(F_\theta\) on any prefix \(U=(u_1,\ldots,u_m)\) of \(d\)-dimensional vectors (tokens or latents). 
Denote the last-layer hidden state at the final position by
\[
H_\theta(U)\in\mathbb{R}^{d}.
\]

For supervision, the \(k\)-th textual step is \(s_k=(y_{k,1},\ldots,y_{k,L_k})\in\mathcal{V}^{L_k}\), and the answer is \(a=(a_1,\ldots,a_{L_a})\in\mathcal{V}^{L_a}\).
The auxiliary decoder has parameters \(\phi\); the LLM has parameters \(\theta\).

\subsection{Implicit Phase: Latent Construction by Last Hidden States}
We fix the number of implicit reasoning steps \(K\) in advance. For each step \(k=1,\ldots,K\),
\begin{equation}
\label{eq:latent}
z_k \;=\; H_\theta\!\big(U^{(k-1)}\big) \in \mathbb{R}^{d}, 
\qquad
U^{(k)} \;=\; U^{(k-1)} \;\oplus\; z_k ,
\end{equation}
where \(\oplus\) denotes concatenation along the time axis. 
The implicit chain-of-thought is therefore represented as a continuous sequence of hidden states \(z_{1:K}=(z_1,\ldots,z_K)\), 
which are autoregressively generated and appended to the context before the model switches to explicit decoding.

\subsection{Explicit Phase: Answer Decoding over the Vocabulary}
After constructing the implicit latents \(z_{1:K}\), the model switches to explicit decoding to generate the final answer. 
Let \(W_o\in\mathbb{R}^{|\mathcal{V}|\times d}\) be the output projection (LM head). 
With teacher forcing on the partial answer \(a_{<t}\), the generation is
\begin{align}
h_{T+K+t} &= H_\theta\!\big(U^{(K)} \oplus e(a_{<t})\big), \\
p_\theta(a_t \mid x, z_{1:K}, a_{<t}) &= \mathrm{softmax}\!\big(W_o\, h_{T+K+t}\big)_{a_t}, \\
p_\theta(a \mid x, z_{1:K}) &= \prod_{t=1}^{L_a} p_\theta(a_t \mid x, z_{1:K}, a_{<t}).
\label{eq:explicit}
\end{align}

\subsection{Training-time Decoder and Step-level Supervision}
During training, a decoder $p_\phi$ (architecturally identical to the LLM) takes only the
$k$-th implicit latent $z_k$ as conditioning signal and autoregressively generates the
$k$-th textual step $s_k=(y_{k,1},\ldots,y_{k,L_k})$. This provides \textbf{step-level}
supervision that directly grounds $z_k$ to its corresponding reasoning content:
\begin{equation}
\label{eq:decoder-factor}
p_\phi(s_{1:K} \mid z_{1:K})
= \prod_{k=1}^{K} p_\phi(s_k \mid z_k)
= \prod_{k=1}^{K} \prod_{t=1}^{L_k}
     p_\phi\!\big(y_{k,t} \mid z_k,\, y_{k,<t}\big).
\end{equation}

\noindent\textit{Parameterization.}
For step $k$, the decoder is conditioned on the implicit latent $z_k$ obtained from the LLM.  
Since $z_k$ does not correspond to any token in the vocabulary, it is not included in the loss calculation.  
Instead, $z_k$ is injected as an additional prefix vector that initializes the decoder’s hidden state for step generation.
Concretely, the decoder input sequence is
\[
U^{\text{dec}}_k = \big[\, z_k \,;\, e(y_{k,1}), \ldots, e(y_{k,L_k}) \,\big],
\]
where $e(\cdot)$ denotes the embedding function of the LLM shared between both models.  
During training with teacher forcing, the decoder predicts each token $y_{k,t}$ autoregressively:
\[
p_\phi(y_{k,t} \mid z_k, y_{k,<t}) = \mathrm{softmax}\!\big(W^{\text{dec}}\, h^{\text{dec}}_{k,t}\big)_{y_{k,t}},
\]
where $h^{\text{dec}}_{k,t}$ is the decoder hidden state at position $t$ and $W^{\text{dec}}$ is the LM head of the decoder.  

The training loss for step $k$ is then
\[
\mathcal{L}_{\text{step},k} = - \sum_{t=1}^{L_k} \log p_\phi(y_{k,t} \mid z_k, y_{k,<t}),
\]
which supervises only the textual step tokens.  
The decoder is used exclusively for this supervision during training and is discarded at inference.

\subsection{Objectives}
Training involves two complementary cross-entropy losses: one for supervising the textual steps through the decoder, and one for supervising the final answer through the base LLM.

\noindent \textbf{Step-level supervision.}
For each implicit latent $z_k$, the decoder $p_\phi$ generates the corresponding reasoning step $s_k=(y_{k,1},\ldots,y_{k,L_k})$.  
Since $z_k$ is not a vocabulary token, the loss is computed only over the textual step tokens:
\begin{equation}
\mathcal{L}_{\text{step}}
= - \sum_{k=1}^{K}\sum_{t=1}^{L_k}
\log p_\phi\!\big(y_{k,t}\mid z_k, y_{k,<t}\big).
\end{equation}
This loss grounds each latent $z_k$ to a specific reasoning step, ensuring that the latent sequence carries fine-grained semantics.

\noindent \textbf{Answer supervision.}
After $K$ implicit steps, the LLM $F_\theta$ switches back to explicit decoding to generate the final answer $a=(a_1,\ldots,a_{L_a})$.  
We optimize the standard language modeling loss:
\begin{equation}
\mathcal{L}_{\text{ans-lm}}
= - \sum_{t=1}^{L_a}
\log p_\theta\!\big(a_t \mid x, z_{1:K}, a_{<t}\big).
\end{equation}

\noindent \textbf{Total objective.}
The overall loss is a weighted sum:
\begin{equation}
\mathcal{L}
= \lambda_{\text{step}}\,\mathcal{L}_{\text{step}}
+ \lambda_{\text{lm}}\,\mathcal{L}_{\text{ans-lm}} .
\label{eq:total}
\end{equation}
Gradients from $\mathcal{L}_{\text{step}}$ propagate through the decoder into the latent representations $z_{1:K}$ and further into the LLM (via Eq.~\eqref{eq:latent}), shaping the hidden states to encode step-level reasoning.  
Meanwhile, $\mathcal{L}_{\text{ans-lm}}$ trains the base model to produce the final answer directly, so the decoder can be discarded at inference time without affecting efficiency.
Implementation details, inference procedures, and diagnostic analyses are provided in Appendix~\ref{app:train-infer-details}.

\section{experiment}
\subsection{Experimental Setup}
\noindent \textbf{Training Data.} We follow previous works \citep{icot, hao2024coconut} to use the \textbf{GSM8k-Aug} dataset \cite{icot} for training implicit CoT models.
The GSM8k-Aug expands the original GSM8k training set \citep{cobbe2021trainingverifierssolvemath} to 385k examples by using GPT-4 for data generation.
To facilitate implicit CoT training, the GSM8k-Aug removes the reasoning chain of natural language, preserving only a sequence of structured mathematical expressions.
Each expression is logically linked to the previous step, as illustrated by the example: \texttt{<<12*3=36>><<9*2=18>><<17*2=34>><<36+18+34=88>>}.

\noindent \textbf{Evaluation Benchmarks.}
We report results on the \textbf{GSM8k-Aug} test set \citep{cobbe2021trainingverifierssolvemath}, which serves as our in-domain (ID) evaluation benchmark.
To further evaluate mathematical reasoning under a distribution shift, we also evaluate models on three out-of-domain (OOD) benchmarks:
(1) \textbf{SVAMP} \citep{patel-etal-2021-nlp}, a dataset of grade-school arithmetic word problems that introduces simple variations to assess robustness;  
(2) \textbf{GSM-Hard} \citep{gao2022pal}, a modified version of the GSM8k test split where numbers are replaced with larger magnitudes to increase problem difficulty; and  
(3) \textbf{MultiArith} \citep{roy2016solvinggeneralarithmeticword}, a subset of MAWPS \citep{koncel-kedziorski-etal-2016-mawps} consisting of multi-step arithmetic word problems.  
Please refer to the Appendix \ref{app:benchmark_detail} for more details.

\input{tables/acc_and_tokens-gpt2-main_table}
\input{tables/acc_and_tokens-llama3_1b-main_table}

\noindent \textbf{Implementation Details.} We follow the training setup of previous works \citep{hao2024coconut,shen2025codi}, and adopt consistent hyperparameter choices for GPT-2, LLaMA 1B/3B/8B. Detailed configurations, such as learning rates, curriculum strategies, are provided in \cref{appendix:implementation}.

\subsection{Main Results}
\noindent \textbf{Baselines.}  
We compare our \methodname against five representative baselines:  
(1) \textbf{CoT-SFT}: Supervised fine-tuning (SFT) on CoT-annotated data, where the model is trained to generate explicit intermediate reasoning steps followed by the final answer.  
(2) \textbf{No-CoT-SFT}: Supervised fine-tuning on direct answers only, without producing intermediate steps.  
(3) \textbf{iCoT} \citep{icot}: A curriculum learning method based on ``Stepwise Internalization,'' which injects CoT reasoning patterns into the model’s internal representations, enabling it to produce more accurate direct answers during inference.  
(4) \textbf{Coconut} \citep{hao2024coconut}: A curriculum learning approach that gradually replaces explicit reasoning steps with implicit tokens until the reasoning process becomes fully implicit. This method has shown strong empirical performance and serves as a primary baseline in our experiments.  
(5) \textbf{CODI} \citep{shen2025codi}: A distillation-based method where explicit CoT acts as the teacher and implicit CoT as the student. By aligning the last hidden states of the full reasoning trajectory, CODI effectively internalizes knowledge and alleviates catastrophic forgetting.

\noindent \textbf{In-Domain Math Benchmark Results.}  
Table~\ref{tab:main_results} (first column) reports GPT-2 results on GSM8k-Aug. \methodname outperforms SFT-CoT and is the first training-based approach where implicit CoT surpasses explicit CoT. With GPT-2 using Coconut as the backbone, it achieves a $+2.1$ point improvement over SFT-CoT. It also exceeds other training-based implicit reasoning models; for example, on Coconut, it improves by $+8.2$ points, a relative gain of $22.4\%$. Moreover, when applied on top of CODI---the current SOTA implicit reasoning method---\methodname yields an additional $+0.6$ point improvement.


Table~\ref{tab:main_results-llama1b} (first column) shows the results when CODI is used as the backbone. In this setting, our method achieves a substantial $+3.4$ point improvement. Furthermore, we are the first to achieve performance comparable to SFT-CoT on LLaMA-1B, reaching $96\%$ of its accuracy. Given that prior studies \citep{xu2025softcot, shen2025codi} reported that curriculum learning in larger models leads to catastrophic forgetting and that homogeneous knowledge harms model training \citep{LLMSurvey}, we choose CODI as the backbone because its KL-regularized objective constrains the training not to deviate too far from the original model distribution, thereby alleviating catastrophic forgetting.

\noindent \textbf{Out-of-Domain Math Benchmark Results.} 
To evaluate the robustness of our method, we train on GSM8k and evaluate on out-of-domain datasets (GSM-Hard, MultiArith, and SVAMP). From the third column of Table~\ref{tab:main_results}, we observe that \methodname consistently outperforms SFT-CoT, with an average improvement of $+4.3$ points when using Coconut as the backbone. From the third column of Table~\ref{tab:main_results-llama1b}, our method further improves upon the current SOTA implicit reasoning method CODI by $+1.0$ point. Moreover, when scaling model size from GPT-2 to LLaMA-1B, \methodname enlarges the performance gap against iCoT, Coconut, and other baselines.  

We attribute the robustness of \methodname to its step-level implicit supervision. Unlike SFT-CoT, which forces the model to mimic deterministic natural language annotations, and unlike CODI, which applies trajectory-level alignment to a coarse-grained reasoning path, our method introduces a moderate form of supervision. This design ensures the plausibility of each reasoning step while preserving the diversity of reasoning trajectories, thereby improving generalization to unseen inputs.

\noindent \textbf{Inference Efficiency.} In terms of inference speed, our method maintains the same efficiency as other implicit reasoning approaches on both GPT-2 and LLaMA-1B. On GPT-2, \methodname not only surpasses SFT-CoT on both in-domain and out-of-domain benchmarks, but also achieves a $2.3\times$ and $2.2\times$ speedup on Coconut, respectively. On LLaMA-1B, \methodname remains comparable to SFT-CoT in accuracy while delivering $1.9\times$ and $1.7\times$ speedups on in-domain and out-of-domain benchmarks, respectively. These results demonstrate the effectiveness of our approach in retaining or even enhancing the performance of explicit CoT while substantially reducing inference cost.

\subsection{Ablation Studies} 
\noindent \textbf{Ablation on the Number of Implicit Tokens.}  
We study the effect of varying the number of implicit latents on GPT-2, comparing \methodname with Coconut trained on GSM8k-Aug and evaluated on GSM8k-Aug, GSM-Hard, MultiArith, and SVAMP (Fig.~\ref{fig:different_latents}). Following Coconut, each latent corresponds to two tokens. As shown in Fig.~\ref{fig:gsm8k-step-distribution}, most problems involve two to six steps with a small proportion of harder cases, so we set the maximum number of implicit latents to 8. For each configuration, we report the best performance, and results show that \methodname provides more stable training and achieves consistent gains over Coconut, indicating that step-level implicit supervision scales effectively with larger latent capacity.
\input{tables/acc_and_tokens-llama3_3b_and_8b}
\input{figures/coconut_simcot_results}

\noindent \textbf{Ablation on Scaling to Larger Backbones.}  
To examine robustness and scalability, we extend experiments to larger LLaMA backbones, including LLaMA 3.2 3B and LLaMA 3.1 8B.  
Table~\ref{tab:llama3-3b_and_8b} reports results on GSM8k-Aug (in-domain) and GSM-Hard, MultiArith, and SVAMP (out-of-domain).  
Overall, \methodname scales effectively to larger backbones, consistently surpassing or matching explicit CoT on out-of-domain tasks while reducing reliance on trajectory-level supervision.  
For clarity, the detailed results and discussion are provided in Appendix~\ref{appendix:scaling_backbones}.

\noindent \textbf{Ablation on Different Decoder Sizes.}  
We replace the decoder of LLaMA 1B with larger ones from the same vocabulary family and evaluate on GSM8k-Aug, GSM-Hard, MultiArith, and SVAMP. A 1B-scale decoder yields consistent gains, whereas larger decoders (3B or 8B) slightly reduce accuracy; detailed discussion is in Appendix~\ref{appendix:decoder_size}.

\noindent \textbf{Ablation on Soft Thinking.} 
We also study the effect of integrating soft thinking \citep{zhang2025soft, wu2025llms} with both Coconut and SIM-CoT. 
For clarity, the detailed experimental setup, results, and analyses are provided in Appendix~\ref{appendix:softthinking}.

% \vspace{-5pt}
% \section{Interpretability Analysis}

\input{figures/Interpretability_analysis_horizontal}

\noindent \textbf{Interpretability of Implicit Reasoning.} 
Continuous thoughts in implicit reasoning models are not aligned with vocabulary tokens and cannot be directly decoded into human-readable text. 
We address this by reusing the training decoder to project and visualize the semantic meaning of each latent step as shown in Fig. \ref{fig:visualization}. 
We also analyze latent token distances under different configurations. 
Detailed descriptions, numerical results, and examples are provided in Appendix~\ref{appendix:interpretability}.

\section{Conclusion}
We introduce \methodname, a training-based implicit reasoning method with step-level supervision on latent tokens. On GPT-2, \methodname outperforms the strong explicit baseline SFT-CoT, while also surpassing implicit baselines such as Coconut and CODI. When scaling to larger LLaMA backbones, the performance achieves consistent gains over existing implicit reasoning methods and maintains fast inference efficiency. Ablation studies further show that it improves training stability with more latent tokens and can benefit from integration with training-free techniques such as soft thinking. Distance analysis confirms that \methodname produces latent representations that are diverse yet stable. 


\appendix
\newpage
\section*{\centering APPENDIX}

\section*{Usage of Large Language Models}
In this paper, we used LLMs only for minor language polishing and formatting, without generating ideas, analyses, or experimental results.

\section*{Outline}
In this appendix, we provide additional analyses and supporting materials to complement the main text.  
First, in Sec.~\ref{appendix:scaling_backbones}, we investigate the scaling behavior of \methodname on larger backbones and show its robustness across different model sizes.  
In Sec.~\ref{appendix:decoder_size}, we analyze the impact of decoder size on performance, highlighting the trade-offs between moderate and excessively large decoders.  
In Sec.~\ref{appendix:softthinking}, we present experiments on combining soft thinking with \methodname, including setup, results, and a detailed formulation with pseudocode.  
In Sec.~\ref{appendix:related}, we provide an overview of related work in explicit and implicit chain-of-thought reasoning.  
In Secs.~\ref{appendix:implementation} and \ref{app:train-infer-details}, we describe our implementation, hyperparameter configurations, and training/inference procedures, including benchmark and dataset details.  
In Sec.~\ref{appendix:simcot}, we introduce the SIM-CoT training procedure and provide pseudocode for step-level supervision.  
In Sec.~\ref{appendix:geometry}, we offer geometric diagnostics of the latent space, analyzing inter-latent distances and distance to the vocabulary center.  
In Sec.~\ref{appendix:interpretability}, we discuss interpretability analysis, including latent visualization, distance analysis, and summary findings.  
Finally, we provide future directions, declarations on LLM usage and additional case studies on GSM8k to further illustrate the reasoning process and visualization choices.


\section{Additional Analysis on Scaling to Larger Backbones}
\label{appendix:scaling_backbones}

As shown in Table~\ref{tab:llama3-3b_and_8b}, \textbf{On LLaMA 3.2 3B}, \methodname improves over CODI by $+1.5$ points on GSM8k-Aug and $+1.6$ points on SVAMP, while maintaining comparable performance on GSM-Hard and MultiArith.  
This indicates that step-level implicit supervision can strengthen strong implicit reasoning baselines.  

\textbf{On LLaMA 3.1 8B}, \methodname yields $+3.0$ points on GSM8k-Aug, $+1.3$ points on SVAMP, and $+0.8$ points on MultiArith relative to CODI, with stable accuracy on GSM-Hard.  
Compared with SFT-CoT, it achieves higher accuracy on MultiArith ($100.0$ vs.\ $98.3$) and SVAMP ($79.4$ vs.\ $73.1$), while remaining similar on GSM-Hard.  

These results confirm that \methodname scales effectively to larger backbones, providing consistent gains while reducing reliance on trajectory-level supervision.

\section{Additional Analysis on Decoder Sizes}
\label{appendix:decoder_size}

Compared to the baseline, as shown in Table~\ref{tab:llama1b-vs-distance}(a), integrating a 1B-scale decoder leads to consistent improvements across all benchmarks.  
However, simply replacing the decoder with larger versions (3B or 8B) does not bring further benefits and slightly degrades performance.


These results suggest that moderate decoder scaling can enhance reasoning ability, but excessively large decoders may introduce optimization difficulties or misalignment with the backbone, ultimately limiting generalization.  
Another plausible explanation is that the 1B encoder and 1B decoder originate from the same model, thereby sharing a more compatible representation space that facilitates learning.  
In contrast, larger decoders (3B or 8B) may require additional projection to align with the 1B backbone, which could introduce mismatches and hinder training stability.


\section{Additional Analysis on Soft Thinking}
\label{appendix:softthinking}

Soft thinking \citep{zhang2025soft, wu2025llms} is a training-free method for implicit reasoning in which the latent space is represented as a weighted average over the vocabulary embedding space.  
In contrast, \methodname learns latent representations directly from data during training.  
To our knowledge, no prior work has evaluated a combination of these two approaches; our experiments provide the first such evaluation.

\subsection{Setup}
We apply the proposed soft thinking mechanism on top of both Coconut and SIM-CoT, while adopting GPT-2 as the backbone model.
To assess the effectiveness of this approach, we perform evaluations on a diverse set of mathematical reasoning benchmarks.
The in-domain evaluation is carried out on GSM8k-Aug, which provides augmented training and testing samples closely aligned with the original GSM8k distribution.
To further examine generalization beyond the training domain, we include three out-of-domain benchmarks: GSM-Hard, which contains more challenging arithmetic problems with subtle variations in reasoning steps; MultiArith, which evaluates performance on multi-step arithmetic operations requiring careful sequencing of addition, subtraction, multiplication, and division; and SVAMP, which focuses on variations of elementary word problems designed to test robustness to superficial changes in problem statements.


\subsection{Results}
Table~\ref{tab:soft-thinking} (b) reports the results.  
Adding soft thinking improves accuracy in most cases.  
For Coconut, improvements are observed on GSM-Hard (+0.2) and MultiArith (+1.7), with a slight decrease on SVAMP (-0.2).  
For SIM-CoT, soft thinking consistently enhances performance: GSM8k-Aug (+0.2), GSM-Hard (+0.1), MultiArith (+0.7), and SVAMP (+0.1).

\subsection{Formulation}
Let $z \in \mathbb{R}^d$ denote a continuous latent token, and $E \in \mathbb{R}^{|\mathcal{V}| \times d}$ be the embedding matrix of the vocabulary $\mathcal{V}$.  
Our goal is to enrich the representational capacity of $z$ by incorporating soft thinking, which allows the latent space to draw information not only from its continuous representation but also from the semantic structure of the vocabulary. The process can be described in three steps.

\paragraph{Step 1. Vocabulary distribution.}  
The continuous latent token $z$ is first mapped into a probability distribution over the vocabulary space:
\[
p = \text{softmax}(W z),
\]
where $W \in \mathbb{R}^{|\mathcal{V}| \times d}$ is the output projection matrix and $p \in \mathbb{R}^{|\mathcal{V}|}$ is the resulting distribution.  
This step can be viewed as interpreting the latent token in terms of vocabulary-level semantics, where each token in $\mathcal{V}$ is assigned a likelihood according to its relevance to $z$.

\paragraph{Step 2. Soft-thinking embedding.}  
Using the distribution $p$, we compute a weighted mixture of vocabulary embeddings:
\[
z_{\text{soft}} = E^\top p = \sum_{v \in \mathcal{V}} p_v \, E_v ,
\]
where $E_v$ is the embedding vector corresponding to token $v$.  
This operation can be seen as constructing a "soft token" that captures multiple semantic hypotheses simultaneously, instead of committing to a single discrete vocabulary token. As a result, $z_{\text{soft}}$ provides richer and smoother information than a hard token lookup.

\paragraph{Step 3. Combination.}  
Finally, we combine the original continuous latent $z$ with the soft-thinking embedding $z_{\text{soft}}$:
\[
z' = \alpha z + \beta z_{\text{soft}},
\]
where $\alpha = \texttt{continuous\_weight}$ and $\beta = \texttt{soft\_weight}$ are hyperparameters that balance the contribution of the continuous and soft-thinking components.  
This formulation allows $z'$ to retain the model’s learned continuous representations while also grounding them in the vocabulary space. Intuitively, the continuous part encourages compact reasoning within the latent space, whereas the soft-thinking component brings in semantic priors from the vocabulary, leading to more stable and interpretable reasoning.  

The pseudocode implementation of the above process is presented as follows.

\begin{algorithm}[h]
\caption{Soft Thinking with Continuous Tokens}
\label{alg:soft-thinking}
\begin{algorithmic}[1]
\Require Continuous latent $z$, embedding matrix $E$, weights $\alpha, \beta$
\If{$\beta > 0$}
    \State Compute logits: $l \gets Wz$
    \State Convert to probabilities: $p \gets \text{softmax}(l)$
    \State Form soft embedding: $z_{\text{soft}} \gets E^\top p$
    \State Update latent: $z' \gets \alpha z + \beta z_{\text{soft}}$
\Else
    \State $z' \gets z$
\EndIf
\State \Return $z'$
\end{algorithmic}
\end{algorithm}


\subsection{Analysis}
The results demonstrate that soft thinking complements training-based implicit reasoning.
The hybrid latent $z'$ integrates semantics learned through training and distributional information from vocabulary mixing, which enables the model to explore diverse intermediate states rather than committing to a single deterministic path.
This leads to improvements in both in-domain and out-of-domain benchmarks.  
Our findings suggest that combining training-free construction with training-based supervision provides gains beyond either approach in isolation.

\input{tables/scale_decoder_and_distance_analysis}
\input{tables/ablation_soft_thinking}

\section{Related Work}
\label{appendix:related}

\noindent \textbf{Explicit chain-of-thought reasoning.}  
Chain-of-thought (CoT) prompting enables large language models (LLMs) to generate intermediate reasoning steps before producing the final answer \citep{wei2023chainofthoughtpromptingelicitsreasoning}.  
This approach has been widely studied and extended in many directions. Self-consistency samples multiple reasoning paths and selects the majority answer to improve reliability \citep{wang2023selfconsistency}.  
Least-to-most prompting decomposes a complex question into simpler sub-problems and solves them in order \citep{zhou2023least}.  
Reflection-based reasoning allows the model to revise or verify its own intermediate steps, leading to better correctness \citep{shinn2023reflexion, madaan2023self}.  
Other works focus on using external tools or symbolic solvers together with explicit reasoning, which further improves accuracy in mathematics and program synthesis \citep{yao2023react}.  
Methods such as progressive-hint prompting \citep{zheng2023progressive} and step-level feedback \citep{wei2025survey} study how supervision can be incorporated into explicit reasoning to make reasoning more structured.  
Despite these advances, explicit CoT has clear drawbacks. Because it generates long token sequences, inference cost grows rapidly with reasoning length, and many intermediate steps are redundant or irrelevant to the final answer.  
Moreover, since explicit reasoning is restricted to tokens from a fixed vocabulary, it often commits to a single trajectory and shows limited reasoning diversity \citep{li2025implicit, zhang2025soft, xu2025softcot}. 

\noindent \textbf{Implicit chain-of-thought reasoning.}
Implicit CoT performs multi-step computation in a continuous latent space instead of emitting long textual traces, reducing decoded length while keeping internal structure. Prior work follows four practical routes. First, \textbf{knowledge internalization} trains models to carry out reasoning internally by progressively removing explicit traces or by using dedicated control embeddings; examples include iCoT-SI \citep{icot}, which removes steps during training to internalize reasoning. Second, \textbf{architectural modification} controls compute by reusing or skipping layers, or by adding light recurrence, so models can refine hidden states without lengthening outputs \citep{loop,innerthinking,ccot,tokenassorted,cotformer,geiping2025scaling}. Third, \textbf{training-free} methods construct continuous latents directly from the model’s probability distribution over the vocabulary; Soft Thinking mixes embeddings by probability to form “concept” tokens that explore alternative paths without updating weights, which improves efficiency and diversity but does not bind each latent to step-level semantics \citep{zhang2025soft,wu2025llms}. 

The fourth route, \textbf{auto-regressive latent reasoning}, updates and concatenates latent states in place of some token-level decoding and is the most relevant to our work \citep{xu2025softcot,tan2025colar}. Coconut applies \textbf{answer-level} supervision—training on the final answer while leaving intermediate latents weakly constrained \citep{hao2024coconut}. CODI adds \textbf{trajectory-level} distillation by aligning an implicit trajectory with an explicit CoT trace, narrowing the gap to explicit CoT but giving only coarse guidance to intermediate steps \citep{shen2025codi}. However, the implicit token length in CODI is fixed during training, which limits its flexibility and makes it less suitable for scaling to variable or longer reasoning chains.
Our framework remains in the auto-regressive setting but changes the supervision: during training, each latent is aligned with its corresponding textual step (\textbf{step-level} supervision), distributing learning signals across the full latent chain to improve stability and semantic fidelity of intermediate states; at inference, the decoder is discarded, ensuring that the decoding cost remains identical to that of standard implicit CoT methods (e.g., Coconut).

\section{Implementation and Training Details}
\label{appendix:implementation}

We provide the full hyperparameter settings, training procedures, and additional analysis used in our experiments. Unless otherwise specified, we use the Adam optimizer with $\beta_1=0.9$, $\beta_2=0.999$, and weight decay of $0.1$. Batch size is set to 128 for GPT-2 and LLaMA 1B, and 64 for LLaMA 3B and 8B. Early stopping is applied with a patience of 3 epochs. We now describe the training setups for Coconut, CODI, and our \methodname.

\subsection{Coconut Training Setup}
Following \citet{hao2024coconut}, GPT-2 and LLaMA 1B \citep{Radford2019LanguageMA,llama3.2} are trained with a fixed learning rate of $1 \times 10^{-4}$. One implicit latent corresponds to two implicit tokens. A curriculum is applied: every three epochs, one explicit reasoning step is replaced by an implicit latent until the maximum number of latent steps is reached. After this expansion, training continues for 15 additional epochs.

\subsection{CODI Training Setup}
For larger backbones such as LLaMA 3B and LLaMA 8B, we adopt task-specific hyperparameter settings to ensure stable training. In particular, we use a learning rate of $3 \times 10^{-4}$ for LLaMA 3B and train for 8 epochs, while for LLaMA 8B the learning rate is reduced to $1 \times 10^{-4}$ with 6 training epochs. These choices are motivated by the increased sensitivity of larger models to optimization dynamics, where smaller learning rates and fewer epochs help to prevent overfitting and instability.  

When reproducing CODI on GPT-2 and LLaMA 1B, we strictly follow the configurations reported by \citet{shen2025codi}. Specifically, we use a learning rate of $3 \times 10^{-3}$ with 40 epochs for GPT-2, and a learning rate of $8 \times 10^{-4}$ with 10 epochs for LLaMA 1B. Adopting these settings ensures that our results are directly comparable to prior work and isolates the effect of our proposed method, rather than confounding it with differences in optimization schedules.  


\subsection{Summary of Hyperparameters}
\begin{table}[h!]
\centering
\caption{Training hyperparameters across different models.}
\label{tab:hyperparams}
% \large
\begin{tabular}{lccc}
\toprule
\textbf{Model} & \textbf{Method} & \textbf{LR} & \textbf{Epochs} \\
\midrule
GPT-2          & Coconut & $1 \times 10^{-4}$ & 15 + curriculum  \\
LLaMA 1B       & Coconut & $1 \times 10^{-4}$ & 15 + curriculum  \\
GPT-2          & CODI    & $3 \times 10^{-3}$ & 40  \\
LLaMA 1B       & CODI    & $8 \times 10^{-4}$ & 10  \\
LLaMA 3B       & CODI    & $3 \times 10^{-4}$ & 8   \\
LLaMA 8B       & CODI    & $1 \times 10^{-4}$ & 6   \\
\bottomrule
\end{tabular}
\end{table}

\section{Training and Inference Details}
\label{app:train-infer-details}

\paragraph{Curriculum for \(K\).}  
We use a curriculum schedule to gradually increase the number of implicit steps. Each latent corresponds to two implicit tokens. Let \(K_{\max}\) denote the maximum number of latents. Starting from \(K^{(0)}=0\), the number of implicit steps after epoch \(e\) is
\[
K^{(e)} \;=\; \min\!\left(K_{\max},\;\Big\lfloor \tfrac{e}{\Delta e} \Big\rfloor \right),
\]
where \(\Delta e\) is the update interval in epochs. Once \(K^{(e)}\) reaches \(K_{\max}\), it remains fixed for the remainder of training.

\paragraph{Inference and Efficiency.}  
At inference time, the auxiliary decoder is removed and only the base model is executed:
\[
U^{(0)} = [\, e(x_1), \ldots, e(x_T) \,], \quad 
\text{for } k=1,\ldots,K:\;
z_k = H_\theta(U^{(k-1)}), \quad
U^{(k)} = U^{(k-1)} \oplus z_k ,
\]
and after $K$ implicit steps the model switches back to explicit decoding to generate the final answer sequence $a$ as in Eq.~\eqref{eq:explicit}.  
The total decoding length is $T + K + L_a$, where $T$ is the input length, $K$ the number of implicit steps, and $L_a$ the answer length. In practice, the cost is comparable to other implicit reasoning methods because $K$ is moderate. In tasks where explicit CoT requires long trajectories ($L_{\text{CoT}} \gg K$), the implicit formulation reduces decoding positions, providing efficiency gains without loss of reasoning accuracy.

\paragraph{Benchmark Detail.}
\label{app:benchmark_detail}
\textbf{SVAMP} \citep{patel-etal-2021-nlp} (Simple Variations on Arithmetic Math Word Problems) is a benchmark dataset designed to test the robustness of math word problem solvers to superficial changes. It contains 1,000 elementary-level arithmetic word problems (grade 4 and below), each involving a single unknown and solvable by an arithmetic expression with no more than two operators. The problems are transformations of existing datasets (such as MAWPS and ASDiv-A) with controlled variations in wording, structure, and number values to reduce artifacts and superficial cues. SVAMP’s average number of reasoning steps required is around 1.2, similar to the base datasets, but model performance drops significantly when tested on SVAMP, showing that many models rely on heuristic patterns rather than deep understanding. 

\textbf{GSM-Hard}~\citep{gao2022pal} is a more challenging variant of the GSM8K dataset, intended to test models’ ability to cope with harder numerical values. It retains the same problem statements as the original GSM8K, but replaces many of the numbers with larger and less common numerical values, making superficial arithmetic computation and reasoning harder. The dataset has about 1,319 examples (matching the GSM8K test set size).

\textbf{MultiArith}~\citep{roy2016solvinggeneralarithmeticword} is a dataset of multi-step arithmetic word problems designed to challenge systems to correctly sequence multiple operations. It contains 600 problems collected from educational sources, each solvable by one equation involving two or more of the four basic operations (addition, subtraction, multiplication, division). The problems require reasoning over multiple sentences to extract and combine numeric quantities, understand the implied operations, and compute the result. MultiArith has been widely used as a benchmark for evaluating arithmetic reasoning generalization, particularly for models that go beyond single-operation problems. Analyses show that many models struggle on these examples compared to simpler datasets, highlighting the importance of handling compositional and sequential numerical reasoning.

\paragraph{Training Data.}
\label{app:training_data}
\input{figures/gsm8k_step_distribution}
\textbf{GSM8K-Aug}~\citep{icot} is the only training corpus we use. It is an augmented dataset derived from GSM8K~\citep{cobbe2021trainingverifierssolvemath}, expanding the original 8.5k training problems to roughly 385k examples through paraphrasing, numerical resampling, and synthetic generation with GPT-4. The distribution of reasoning steps in GSM8K-Aug is illustrated in Fig.~\ref{fig:gsm8k-step-distribution}, where the majority of problems require two to four steps, while a long-tail of six or more steps persists. This balance of common and complex instances makes GSM8K-Aug particularly suitable for training models that need to generalize across reasoning difficulty levels.

\section{SIM-CoT Training Implementation}
\label{appendix:simcot}

We provide pseudocode for the SIM-CoT training process, which illustrates how continuous latent embeddings are aligned with explicit supervision at the step level. In particular, each reasoning step in the explicit chain is mapped to a corresponding latent representation, and the training objective enforces consistency between the predicted latent tokens and the ground-truth step annotations. This design ensures that the model learns to represent intermediate reasoning steps in a compact latent space while still retaining interpretability through explicit alignment. By supervising at the step level rather than only at the final answer or trajectory level, SIM-CoT enables finer control over the reasoning process and reduces instability that often arises when scaling to longer chains or larger numbers of implicit tokens.


\begin{algorithm}[h]
\caption{SIM-CoT Training Procedure}
\label{alg:simcot}
\begin{algorithmic}[1]
\Require Batch size $b$, number of thoughts $C$, continuous embeddings $Z$, tokenized inputs $X$, embedding matrix $E$
\For{each thought $t = 1, \ldots, C$}
    \For{each sample $i = 1, \ldots, b$}
        \State Extract continuous embeddings $z_{i,t}$ from $Z$
        \State Obtain token embeddings $e_{i,t}$ from $E(X_{i,t})$
        \State Concatenate embeddings: $h_{i,t} \gets [z_{i,t}; e_{i,t}]$
        \State Build attention mask $m_{i,t}$ up to EOS
        \State Assign position ids $p_{i,t}$
        \State Prepare labels $y_{i,t}$ with masked tokens set to $-100$
    \EndFor
\EndFor
\State Pad and stack $\{h,m,p,y\}$ to maximum sequence length
\State Prepare 4D attention mask: $\hat{M} \gets \text{PrepareMask}(M)$
\State Forward pass: $\hat{O} \gets \text{ExplainableLLM}(H, \hat{M}, P)$
\State Extract logits: $L \gets \hat{O}.\text{logits}$
\State Shift logits and labels: $L' \gets L[:, :-1], \; Y' \gets Y[:, 1:]$
\State Compute cross-entropy loss: $\ell = \text{CrossEntropy}(L', Y')$
\State Normalize $\ell$ over valid positions
\Ensure Final training loss $\ell$
\end{algorithmic}
\end{algorithm}

\section{Geometric Diagnostics of the Latent Space}
\label{appendix:geometry}

We analyze the geometry of latent representations with two metrics.

\textit{Inter-latent distance.}  
\begin{equation}
\mathrm{Dist}(z_{1:K})
= \frac{2}{K(K-1)} \sum_{1 \le i < j \le K} \big\| z_i - z_j \big\|_2 .
\end{equation}
A larger value indicates better separation, reducing the risk of collapse.

\textit{Distance to vocabulary center.}  
Let $\mu = \tfrac{1}{|\mathcal{V}|}\sum_{v\in\mathcal{V}} E_v$ denote the mean embedding. Then,
\begin{equation}
\mathrm{DistVC}(z_{1:K})
= \frac{1}{K} \sum_{k=1}^{K} \big\| z_k - \mu \big\|_2 .
\end{equation}
Moderate values indicate that latents remain close enough to the lexical manifold for stability, while avoiding collapse toward the center. These diagnostics are not used in training but serve as indicators of diversity and stability in the learned latent space.

\input{figures/distance_visulize}

\section{Additional Details for Interpretability Analysis}
\label{appendix:interpretability}

\subsection{Making Implicit Reasoning Visible}
Continuous thoughts produced by implicit reasoning models are represented as latent embeddings that do not correspond to discrete vocabulary tokens, and therefore cannot be directly decoded by a tokenizer. This makes it difficult to interpret how the model internally organizes multi-step reasoning. To address this, we reuse the decoder that was employed for step-level supervision during training, and apply it at inference time to map each latent embedding into a human-readable token sequence.  

As illustrated in Figure~\ref{fig:visualization}, the process begins with a natural language problem (e.g., a math word problem) that is embedded and passed into the large language model. The model generates a sequence of implicit latent tokens, which capture intermediate reasoning steps in continuous space. These latent tokens are then fed into the optional decoder, which translates them into interpretable expressions. Each latent corresponds to one reasoning step, and the autoregressive generation order encodes the dependency structure across steps.  

For example, in the GSM8k case study shown in the figure, the first latent is decoded as $0.3 \times 120 = 36$, representing the number of watermelons harvested initially. The second latent builds upon this result to compute $120 - 36 = 84$, the remaining melons. The third latent then calculates $\tfrac{3}{4} \times 84 = 63$, and the final latent derives the answer $84 - 63 = 21$. For clarity and to save space, the figure merges the second and third steps into a single box, but the actual implicit reasoning unfolds across four distinct latent steps. This sequence of decoded latents mirrors the logic of explicit chain-of-thought reasoning, while being produced implicitly within the latent space.
  

By projecting implicit tokens into interpretable space, we gain direct visibility into how the model structures multi-step reasoning. This not only enables analysis of the correctness and consistency of intermediate steps, but also highlights the dependencies across steps that underlie the final prediction. The visualization confirms that SIM-CoT can encode semantically meaningful and logically ordered reasoning steps in its latent space, bridging the gap between implicit and explicit reasoning.


\subsection{Distance Analysis}
We analyze the geometry of latent representations by measuring two quantities: the average pairwise distance among latent tokens (Dist.) and the average distance from latent tokens to the vocabulary center (Dist. to VC).  
Table~\ref{tab:llama1b-vs-distance}(b) reports the numerical outcomes under different configurations, while Figure~\ref{fig:distance_visualize} provides a visualization of the same phenomena.

As the number of latent tokens increases from 1 to 5, Dist. gradually grows from 20.30 to 28.34, suggesting that the model distributes the latent representations more sparsely in the embedding space, which improves separability. However, in the failed case with 5 latents, Dist. collapses to 4.21, reflecting a degeneration where all latent tokens converge to nearly identical points, losing their ability to represent distinct reasoning steps. By contrast, SIM-CoT pushes Dist. to 32.81 under the same configuration, showing that step-level supervision effectively enforces stronger separation and prevents collapse.

For Dist. to VC, we observe the highest value with 1 latent (36.20). As the number of latents increases, Dist. to VC decreases toward approximately 28, meaning that the representations become better aligned with the vocabulary manifold rather than drifting away. The failed 5-latent case instead shows an abnormal increase to 39.39, a clear signal of instability where the model drifts away from the vocabulary space. SIM-CoT stabilizes this measure at 29.80, achieving a balance: the latent tokens remain distinct from each other while staying anchored close enough to the vocabulary space to preserve semantic interpretability.

The visualization in Figure~\ref{fig:distance_visualize} corroborates these trends. Normal implicit tokens (first row) are well separated and positioned near the vocabulary manifold. Failed tokens (second row) collapse together and drift outward, losing both separation and grounding. After applying SIM-CoT (third row), the tokens regain a structured configuration: distances between latents expand, while their proximity to the vocabulary center remains stable. This alignment between quantitative and qualitative results highlights the effectiveness of SIM-CoT in maintaining a structured and interpretable latent space.

\subsection{Summary}
Overall, the results demonstrate that SIM-CoT establishes a balance between \textbf{diversity} and \textbf{stability} in the latent space. Larger inter-latent distances mitigate representation collapse, while moderate distances to the vocabulary center prevent excessive drift. This equilibrium supports stable implicit reasoning and provides robustness when scaling to more latent tokens.

\section{Future Directions}

We outline several promising directions for extending this work:  
\textbf{(1) Multimodal extension.} Incorporating intermediate supervision from images \citep{chen2023sharegpt4v, sun2024alpha, liu2024mmdu} or videos \citep{chen2024sharegpt4video, wei2025videorope} could generalize implicit CoT to multimodal reasoning tasks \citep{shen2025efficient, zhang2025booststep}, an especially timely direction given the rapid progress of multimodal learning \citep{liu2023llava, chen2023sharegpt4v, zhang2025sec, liu2025songgen, wei2025videoropepp, zhang2024internlm, dong2024internlm, zhao2025omnialign, ding2025mmifenginemultimodalinstructionfollowing, chen2024open, xing2025scalecap, xing2024pyramiddrop}. 
\textbf{(2) Multi-path implicit reasoning.} Exploring branched latent structures, inspired by Tree-of-Thought \citep{yao2023tree, wang2023selfconsistency, zhang2023evaluating} methods, may enhance the diversity and robustness of implicit reasoning.  
\textbf{(3) Integration with RLHF.} Step-level supervision can complement preference optimization and reinforcement learning \citep{guo2025deepseek, rafailov2023direct, liu2025visual, internlmxcomposer2_5_reward}, leading to a more reliable and adaptive reasoning framework.  
\textbf{(4) Theoretical foundations.} Future work may also study implicit reasoning from the perspectives of information theory and representation learning \citep{zhu2025reasoning, xu2025cot}, providing a principled explanation for why step-level supervision stabilizes the latent space.



\section{Additional Case Studies on GSM8k}
In practice, implicit reasoning continues to produce latent tokens even after the correct answer has been reached. These trailing latents no longer introduce new steps but simply repeat the final prediction. For clarity, we omit such redundant tokens in the visualizations. As a result, only the latents that correspond to meaningful intermediate steps are displayed in Figure~\ref{fig:case_all}, while those mapping directly to the final answer are hidden. This choice improves readability without changing the underlying reasoning process.

Notably, the decoded reasoning steps consistently match the semantic structure of explicit chain-of-thought annotations, while being generated implicitly within the latent space. The final predictions align with the ground-truth answers, demonstrating that \methodname is capable of encoding interpretable and step-ordered reasoning without requiring explicit supervision at inference time.  

\input{figures/case_all}


